% Section 4.4: CoT 增强方案设计

CoT 分支把「推理」从隐式激活变为显式两阶段流程，对应 4.1 节矩阵中 H2、H3、H5、H6、H8、H10 六组评测。相对 Direct，CoT 的方法差异集中在推理侧 prompt 结构与生成调度——训练侧由 4.3 节的 LoRA 方案统一承接，本节不再重复。本节 4.4.1 给出 Zero-shot CoT 的两阶段流水线骨架，4.4.2 说明 Few-shot 示例的检索与拼接方式，4.4.3 把 6 组 CoT 对照按设问对齐到矩阵中具体实验编号，4.4.4 汇总为保证两阶段推理稳定而加入的关键工程约束。

\subsection{Zero-shot CoT 流程}
% Subsection 4.4.1: Zero-shot CoT 流程

Zero-shot CoT（CoT-ZS）对应 H2/H5/H8 三组推理侧。其核心思想是显式切分「推理」与「编码」两步，各自用专属 system 指令约束输出形态，图~\ref{fig:ch4-cot-pipeline} 给出整体数据流。

\begin{figure}[!htb]
  \centering
  \begin{tikzpicture}[
    node distance=5mm and 6mm,
    box/.style={draw, rounded corners=2pt, align=center, minimum height=8mm, minimum width=26mm, font=\footnotesize},
    lbl/.style={font=\scriptsize\itshape, align=center},
    arr/.style={-{Stealth[length=2mm]}, thick},
  ]
    \node[box] (q) {题目 $q$};
    \node[box, right=of q] (s1) {Stage-1\\ Reasoning LLM};
    \node[box, right=of s1] (r) {\texttt{reasoning}};
    \node[box, right=of r] (s2) {Stage-2\\ Code LLM};
    \node[box, right=of s2] (c) {\texttt{solution}};

    \draw[arr] (q) -- (s1);
    \draw[arr] (s1) -- (r);
    \draw[arr] (r) -- (s2);
    \draw[arr] (q.south) to [bend right=14] (s2.south);
    \draw[arr] (s2) -- (c);

    \node[lbl, below=1mm of s1] {SYSTEM: algorithm teacher};
    \node[lbl, below=1mm of s2] {SYSTEM: Python programmer};
    \node[lbl, below=1mm of r]  {截断 + 长度预算};
  \end{tikzpicture}
  \caption{Zero-shot CoT 两阶段推理数据流}
  \label{fig:ch4-cot-pipeline}
\end{figure}

\paragraph{（1）Stage-1：推理生成}
第一阶段只负责产出自然语言推理，不允许直接给出代码：System 为「\texttt{expert algorithm teacher. Generate step-by-step reasoning.}」；User 采用与 \texttt{LoRA\_cot} 训练相同的 CoT 版题目指令（题面 + 末尾 \texttt{Think step by step, then output the Python code.}），\textbf{不是} 4.2.1 节 Direct 模板中的 \texttt{Only output the Python code, no explanation.}，以避免「要求只输出代码」与「本阶段只生成推理」相矛盾；采样参数 \texttt{MAX\_TOKENS\_REASONING=1024}、\texttt{temperature=0.3}、\texttt{top\_p=0.95}、\texttt{top\_k=50}，与 Direct 一致。输出清洗：若模型提前滑出任意三反引号代码块（\verb|```python|、\verb|```bash| 等任意语言标记或无标记围栏），立刻截断只保留前置自然语言段，避免「推理阶段偷跑代码」污染第二阶段上下文。逐字模板见附录~\ref{appendix:prompts}~A.3。

\paragraph{（2）推理段的长度预算}
第一阶段输出不能无限拼进第二阶段，否则会压缩 Stage-2 的生成预算、并在 CoT-FS 模式下挤出示例内容。为此引入两级截断：
\begin{itemize}
    \item \textbf{字符级}：\texttt{MAX\_REASONING\_CHARS=1200}，按字符截断长推理；
    \item \textbf{token 级}：\texttt{MAX\_REASONING\_PROMPT\_TOKENS=256}，按 tokenizer 度量再做一次硬截断，保证即使中文/符号密集也能控制在 Stage-2 prompt 中的占位；
    \item \textbf{VLLM 上下文}：\texttt{VLLM\_MAX\_MODEL\_LEN=8192}，相比 Base 的 2048 显著放宽，为两阶段上下文与 few-shot 示例留出空间。
\end{itemize}

\paragraph{（3）Stage-2：代码生成}
第二阶段接收 \{题目, 截断后的推理\} 作为 chat 上下文，要求输出可直接执行的代码：System 为「\texttt{expert Python programmer. Based on the reasoning, write the code.}」；Assistant 期望格式固定「推理文字 + 空行 + \verb|```python ... ```|」，与训练阶段 assistant 回复格式严格一致；采样参数 \texttt{MAX\_TOKENS\_CODE=1024}、其余与 Stage-1 相同。代码抽取：\texttt{eval\_lora\_cot\_passk.py} 中取\textbf{最后一个} \verb|```python| 块（若无则取最后一个无语言围栏块，再退回全文），以便在模型输出多段围栏时仍以最终代码块为准；此逻辑与 4.2.2 节 Direct 路径「取首段围栏」不同，两份 \texttt{extract\_code} 实现仅在 \texttt{[0]} / \texttt{[-1]} 上有差异。逐字 Stage-2 综合 User 模板见附录~\ref{appendix:prompts}~A.3。

\paragraph{（4）端到端评测入口}
CoT-ZS 的完整评测由 \texttt{evaluation/eval\_lora\_cot\_passk.py} 在 \texttt{--few-shot-k 0} 下启动；数据加载、子进程隔离执行、pass@k 无偏估计与 4.2.2 节\textbf{口径}一致，但\textbf{代码抽取规则}见 4.4.1（3）末条（取末段围栏）。实现上将「一次生成 $N{\times}10$ 份结果」拆成「先批量生成 reasoning、再批量生成 code」两次 \texttt{llm.generate} 调用，以匹配两阶段流水线。

\subsection{Few-shot CoT 注入}
% Subsection 4.4.2: Few-shot CoT 注入

Few-shot CoT（CoT-FS($k$)）在 CoT-ZS 基础上的要点如下：
\begin{itemize}
    \item \textbf{示例形态}：每条为 \texttt{(题目, 推理, 代码)}，在 chat 中编码为 \textbf{User(CoT 版题指令)} $\rightarrow$ \textbf{Assistant(推理 + fenced Python)}；
    \item \textbf{两阶段对称回放}：Stage-1 在算法教师 system 之后插入 $k$ 组示范，再接目标题 User；Stage-2 在程序员 system 之后用\textbf{同一组}示例以\textbf{相同顺序}再回放一遍，最后接 \texttt{Reasoning:} + 截断 Stage-1 + 目标题的综合 User；
    \item \textbf{实现对应}：\texttt{eval\_lora\_cot\_passk.py}；训练侧 few-shot 交替结构见 \texttt{train\_cot.py}；
    \item \textbf{实验编号}：H3、H6、H10；
    \item \textbf{选池原则}：按强标签检索，训练与评测共享同一评分函数，以保证示例与目标题尽量同构。
\end{itemize}

\subsubsection{候选评分与选池}
对每条目标样本 $q$，从 3.3.2 节定义的强标签集合派生 \texttt{strong\_tags}，在候选集合（训练集或评测集扣除 $q$）上按下式打分：
\[
\mathrm{score}(q, q') = 2 \cdot \bigl|\,\mathrm{strong\_tags}(q) \cap \mathrm{strong\_tags}(q')\,\bigr|
                       + \bigl|\,\mathrm{tags}(q) \cap \mathrm{tags}(q')\,\bigr|.
\]
\begin{itemize}
    \item \textbf{加权含义}：强标签命中权重 2，普通标签权重 1，兼顾「算法骨架一致」与「题型相近」；
    \item \textbf{候选池规模}：按分数降序取前 $\max(k, 4k)$ 条；
    \item \textbf{池内抽样}：再随机抽样，避免总是固定命中极少数高分题。
\end{itemize}

\subsubsection{训练/评测两侧的随机性差异}
\begin{itemize}
    \item \textbf{共同点}：两侧共用同一评分函数与同一强标签集合，差异仅在随机采样；
    \item \textbf{训练侧}：\texttt{random.sample(pool, k)}，依赖 \texttt{seed=42} 的全局 RNG，每个 epoch 的示例在统计上来自同一分布；
    \item \textbf{评测侧}：\texttt{seed = 2026 + problem\_idx * 1000 + sample\_idx}，同一 \texttt{(题目, 采样编号)} 的示例选择完全可复现，便于跨实验、跨运行做差分归因；
    \item \textbf{合取含义}：训练与评测在「评分 + 候选池构造」上统计一致；评测侧因固定 seed，任意第三方可按题号与采样编号逐条复核。
\end{itemize}

\subsubsection{示例拼接与预算分配}
\begin{itemize}
    \item \textbf{多轮形态}：$k$ 条示例在 Stage-1 与 Stage-2 中均为「\textbf{User（CoT 版题目指令）} $\rightarrow$ \textbf{Assistant（推理 + fenced Python）}」，置于目标轮次之前（与 4.1 节（1b）一致），再接目标题 User 或 \texttt{Reasoning:} 综合 User；
    \item \textbf{三类长度预算}（两阶段 prompt 共享）：
    \begin{itemize}
        \item 示例推理保持原文，不在此阶段额外截断（长度由数据生成阶段约束）；
        \item 目标题的 Stage-1 输出仍受 \texttt{MAX\_REASONING\_CHARS} 与 \texttt{MAX\_REASONING\_PROMPT\_TOKENS} 二级截断；
        \item \texttt{VLLM\_MAX\_MODEL\_LEN=8192} 覆盖「全部示例 + 目标题 + 两阶段生成」。
    \end{itemize}
\end{itemize}

\subsubsection{训练/推理 $k$ 匹配的强约束}
\begin{itemize}
    \item \textbf{规则来源}：4.1 节合法性约束；训练侧为 \texttt{LoRA\_cot\_fs($k$)} 时，推理 \texttt{--few-shot-k} 必须同为 $k$；
    \item \textbf{脚本行为}：不匹配时评测脚本显式退出（H10 对应 \texttt{FS\_K}）；
    \item \textbf{归因意义}：保证「训练见 $k$ 例、评测亦见 $k$ 例」，是 few-shot 净效应可比的前提；否则 train/test 在上下文长度与示例密度上错位，会混淆训练侧与推理侧贡献。
\end{itemize}

\subsection{CoT 对照实验设计}
% Subsection 4.4.3: CoT 对照实验设计

本节与前后文的衔接关系如下：
\begin{itemize}
    \item \textbf{4.4.1--4.4.2}：已给出 CoT-ZS / CoT-FS 的流程与选池；
    \item \textbf{4.1 节（1a）（1b）}：集中给出各组评测时模型见到的 messages 字面骨架（含 Direct 与 CoT 训练侧对照），便于将 H1--H10 与「输入长什么样」逐组对应；
    \item \textbf{本节任务}：将 H2、H3、H5、H6、H8、H10 共 6 组 CoT 评测按设问重组为三块对照，作为第五章归因的直接依据。
\end{itemize}

\subsubsection{推理侧主效应（训练侧固定为 Base）}
以 H1 为共同基线，抽取不依赖任何 LoRA 微调的 CoT 推理侧贡献：
\begin{itemize}
    \item \textbf{$H2 - H1$}：CoT-ZS 相对 Direct 的主效应，回答「零示例推理是否在基座上就有增益」；
    \item \textbf{$H3 - H1$}：CoT-FS($k$) 相对 Direct 的主效应，回答「加 $k$ 条示例是否进一步改进基座」；
    \item \textbf{$H3 - H2$}：few-shot 相对 zero-shot 的边际贡献，隔离示例本身的增益；
    \item \textbf{对照解释}：三条差分共享同一训练侧（Base），仅推理侧变化，故组内显著差异可归因于 prompt 结构。
\end{itemize}

\subsubsection{训练 $\times$ 推理交互效应}
固定 LoRA 作为训练侧、CoT 作为推理侧，通过差分差分量（difference-in-differences）判断两者是否正交可加：
\begin{itemize}
    \item \textbf{$(H8 - H7) - (H2 - H1)$}：在 \texttt{LoRA\_cot\_zs} 与 Base 上「加 CoT-ZS 推理」两条边际是否一致；
    \item \textbf{$(H10 - H9) - (H3 - H1)$}：在 \texttt{LoRA\_cot\_fs($k$)} 与 Base 上「加 CoT-FS($k$) 推理」两条边际是否一致（$k$ 匹配约束下）；
    \item \textbf{$(H5 - H4) - (H2 - H1)$}：在 \texttt{LoRA\_code} 与 Base 上「加 CoT-ZS 推理」两条边际是否一致，额外控制了「未监督 CoT 格式」的训练侧；
    \item \textbf{$(H6 - H4) - (H3 - H1)$}：同上，推理侧换为 CoT-FS($k$)；
    \item \textbf{对照解释}：直接对应 4.1 节「LoRA 与 CoT 是否正交可加」的设问，符号与量级即析因设计的主落点。
\end{itemize}

\subsubsection{Few-shot 边际在训练侧上的一致性}
只变化训练侧，观察 few-shot 的边际贡献是否稳定：
\begin{itemize}
    \item \textbf{$H3 - H2$（Base）vs $H6 - H5$（LoRA\_code）}：两条「$k$ 例相对 $0$ 例」的边际在\textbf{未显式监督 CoT 格式}的训练侧上是否一致；
    \item \textbf{$H10 - H8$（匹配 $k$ 的 LoRA\_cot）}：训练与推理 $k$ 同步放大时，CoT 的净收益相对纯 CoT-ZS 路径的提升量；
    \item \textbf{设问}：few-shot 是否为与训练侧无关的「结构性增益」；
    \item \textbf{读法}：若两条边际量级接近，则 few-shot 主要作用于推理侧；若差异显著，则训练侧调节其有效性，结论中需分情况讨论。
\end{itemize}

\subsubsection{对照汇总}
\begin{itemize}
    \item 三组设计均为「单因子变动、其余固定」，覆盖全部与 CoT 相关的显著差分读法；
    \item 结果文件命名遵循表~\ref{tab:ch4-matrix} 末列；
    \item 第五章仅需按表取数作差，无需补做实验。
\end{itemize}

\subsection{关键实现细节与合法性约束}
\input{chapters/ch04-method/sections/sec04-cot/subsections/04-implementation-constraints}